\chapter{Modelos de Mudança de Regime}
\label{cap:regime}
% ── índice ──────────────────────────────────────────────────────
\index{Markov Switching|textbf}
\index{mudança de regime|textbf}
\index{markov@\texttt{markov()}|textbf}
\index{Hamilton, modelo de|textbf}
\index{probabilidade de transição|textbf}
\index{filtragem de Hamilton}
\index{duração esperada do regime}

% ─────────────────────────────────────────────────────────────────────────────
\section{Heterogeneidade estrutural ao longo do tempo}

Muitas séries econômicas exibem mudanças estruturais: diferentes médias,
volatilidades ou dinâmicas em distintos períodos. Dois enfoques capturam
essa heterogeneidade:

\begin{itemize}
  \item \textbf{Markov Switching} (Hamilton 1989): regime latente segue uma
  cadeia de Markov. A probabilidade de transição entre regimes é estimada
  pelo algoritmo EM junto com os parâmetros de cada regime.
  \item \textbf{Threshold Regression} (Hansen 1997): mudança de regime é
  determinada por uma variável observável (variável limiar $q$) cruzar um
  limiar $\gamma$ estimado.
\end{itemize}

\subsection*{Markov-Switching AR(p)}

\[
  y_t = \mu_{S_t} + \sum_{k=1}^p \phi_k^{(S_t)}\,(y_{t-k} - \mu_{S_t}) + \sigma_{S_t}\,\varepsilon_t
\]

onde $S_t \in \{1,\ldots,K\}$ é o regime latente com probabilidades de
transição $P(S_t = j \mid S_{t-1} = i) = p_{ij}$. O parâmetro $\mu_{S_t}$
é o intercepto do regime, $\phi_k^{(S_t)}$ são os coeficientes AR e
$\sigma_{S_t}$ é o desvio-padrão do regime.

% ─────────────────────────────────────────────────────────────────────────────
\section{Verificação por DGP}

DGP com dois regimes alternantes:
\begin{itemize}
  \item Regime 1 (expansão): $\mu_1 = 1{,}0$, $\phi_1 = 0{,}30$, $\sigma_1 = 0{,}5$
  \item Regime 2 (recessão): $\mu_2 = -1{,}0$, $\phi_2 = 0{,}70$, $\sigma_2 = 1{,}5$
\end{itemize}

\begin{lstlisting}[caption={DGP Markov Switching AR(1) e estimação}]
seed(42)
generate df t      = _n
generate df regime = (t % 60 < 40)   // ~2/3 no regime 1
// y recursivo com parâmetros condicionais ao regime
...
let m_ms = markov(df, y, k=2, p=1)
display m_ms
\end{lstlisting}

\subsection*{Markov Switching}

\begin{lstlisting}[language={}, basicstyle=\ttfamily\small,
  frame=none, numbers=none, backgroundcolor=\color{codbg},
  xleftmargin=0pt, xrightmargin=0pt, breaklines=false]
===================== Markov Switching AR(1) — 2 regimes =====================
Observations: 299   Log-L: -155.76   AIC: 327.51   BIC: 357.12
--------------------------- Transition Matrix --------------------------------
[ 0.9493   0.0507 ]
[ 0.0251   0.9749 ]
--------------------------- Regime 0 (Recessão) ------------------------------
Intercept: -1.2864   AR.L1: 0.6169   Variance: 0.5141
--------------------------- Regime 1 (Expansão) ------------------------------
Intercept:  0.9946   AR.L1: 0.2956   Variance: 0.0625
--------------------------- Expected Durations --------------------------------
Regime 0: 19.7 períodos    Regime 1: 39.9 períodos
\end{lstlisting}

O modelo recupera os dois regimes com boa precisão:
\begin{itemize}
  \item Regime 0 (recessão): $\hat\mu = -1{,}29$ (verdadeiro: $-1{,}0$),
  $\hat\phi = 0{,}62$ (verdadeiro: 0{,}7). Duração esperada 19{,}7 períodos.
  \item Regime 1 (expansão): $\hat\mu = 0{,}99 \approx 1{,}0$,
  $\hat\phi = 0{,}30$ (verdadeiro: 0{,}3). Duração esperada 39{,}9 períodos.
\end{itemize}

A matriz de transição indica alta persistência de ambos os regimes
($p_{00} = 0{,}95$, $p_{11} = 0{,}97$) — a série raramente troca de regime,
condizente com o DGP ($60\%$ do tempo no regime 1).

% ─────────────────────────────────────────────────────────────────────────────
\section{Sintaxe}

\begin{lstlisting}
markov(df, y, k=2, p=1)    // k regimes, AR de ordem p
\end{lstlisting}

\begin{center}
\begin{tabular}{lll}
\toprule
\textbf{Opção} & \textbf{Padrão} & \textbf{Descrição} \\
\midrule
\texttt{k} & 2 & número de regimes \\
\texttt{p} & 1 & ordem AR dentro de cada regime \\
\bottomrule
\end{tabular}
\end{center}

\begin{nota}
  O número de regimes $k$ é fixo e deve ser especificado a priori.
  Critérios de informação (AIC, BIC) comparados entre modelos com $k=2$ e
  $k=3$ guiam essa escolha. Para dados com $N < 200$, $k > 3$ raramente é
  identificado de forma confiável.
\end{nota}
